\chapter{Objetivos}
\label{cap:obj} % etiqueta para poder referenciar luego en el texto con ~\ref{sec:intro}
En este capítulo se exponen el objetivo general de este Trabajo Fin de Grado y los objetivos específicos que se han desarrollado para alcanzar dicho objetivo general. Asimismo se describe la planificación temporal del proyecto, detallando las diferentes fases de desarrollo y las tareas que se han llevado a cabo en cada una de ellas.


\section{Objetivo General}
\label{sec:objgen}

\section{Objetivos Específicos}
\label{sec:objesp}
Con el fin de alcanzar el objetivo general del Trabajo se han desarrollado los siguientes objetivos específicos para facilitar el desarrollo del mismo:
\begin{itemize}
    \item Analizar el funcionamiento e implementación de la versión original de \textit{Dr. Snap!} para identificar y desarrollar o resolver posibles áreas de mejora y problemas de funcionamiento.
    \item Implementar un \textit{Batch Mode}, que permita analizar gran cantidad de proyectos de \textit{Snap!} de forma simultánea y devuelva los resultados en un formato que facilite su análisis individual y conjunto.
    \item Blindar \textit{Dr. Snap!} frente a posibles amenazas para proteger el servidor en el que se despliegue.
    \item Encontrar una ubicación y desplegar una \textit{versión beta} de \textit{Dr. Snap!} abierta a propuestas de mejora.
    \item Validar el correcto funcionamiento y aportación al desarrollo del pensamiento computacional de las mejoras y nuevas funcionalidades incluidas en \textit{Dr. Snap!} a través de la realización de pruebas experimentales y el análisis de los resultados obtenidos en estas.
\end{itemize}

\section{Planificación Temporal}
\label{sec:plan}
El desarrollo de este Trabajo Fin de Grado se ha extendido durante aproximadamente diez meses; en concreto de noviembre de 2025 a septiembre de 2026. Durante este periodo, se han dedicado aproximadamente 3 horas diarias durante los días laborables en horario no lectivo al desarrollo del presente proyecto, existiendo cierta variabilidad en la cantidad de horas diarias dependiendo de variables externas.

Teniendo en cuenta lo anterior, el proyecto se ha desarrollado a través de la consecución de diferentes fases que denominaremos \textit{Sprints}. Cada \textit{Sprint} incluye tareas específicas necesarias para el alcance de los objetivos con anterioridad descritos. A continuación se describen brevemente estas fases:
\begin{itemize}
    \item \textbf{Sprint 1 (noviembre 2025 - diciembre 2025)} Estudio y análisis de la versión de \textit{Dr. Snap!} desarrollada por Del Prado y búsqueda de problemas de funcionamiento y áreas de mejora. Para ello, se realizó una revisión del código fuente\footnote{El código fuente de la versión de Del Prado de \textit{Dr. Snap!} se puede encontrar en el siguiente repositorio de GitHub: \url{https://github.com/mdelprado2017/Dr_Snap}}, así como un estudio del funcionamiento de \textit{Snap!} y comparativa de las funcionalidades con \textit{Dr. Scratch}\footnote{Podemos acceder a la aplicación web a través de la siguiente dirección: \url{https://www.drscratch.org/}}. 
    \item \textbf{Sprint 2 (diciembre 2025 - enero 2026)} Implementación del análisis de un proyecto por URL a través de la API de \textit{Snap!} y resolución de problemas de traducción.
    \item \textbf{Sprint 3 (enero 2026 - febrero 2026)} Desarrollo del \textit{Batch Mode} y refactorización del código para mejorar la eficiencia y legibilidad del mismo. Además de modificar la lógica de parseo de los proyectos \texttt{.xml}.
    \item \textbf{Sprint 4 (Primera mitad febrero 2026)} Adición del análisis de \textit{scripts} duplicados dentro de los \textit{Malos Hábitos} y resolución de problemas encontrados en el análisis de código muerto.
    \item \textbf{Sprint 5 (Segunda mitad febrero 2026)} Modificación del logotipo de \textit{Dr. Snap!} y realización de cambios en la interfaz de usuario para mejorar la experiencia de usuario y ajustar al nuevo logotipo.
    \item \textbf{Sprint 6 (marzo 2026 - abril 2026)} Mejora de la ciberseguridad de la aplicación web, realizar tests de funcionamiento de las medidas aplicadas. Asimismo, inclusión de un aviso de que es una \textit{versión beta} y un formulario de contacto para que los usuarios puedan enviar propuestas de mejora.
    \item \textbf{Sprint 7 (abril 2026 - mayo 2026)} Despliegue de la \textit{versión beta} de \textit{Dr. Snap!} en un servidor y realización de pruebas de funcionamiento y de ciberseguridad.
    \item \textbf{Sprint 8 (mayo 2026 - julio 2026)} Refactorización del código y plantillas, así como, eliminación de los que ya no sean necesarios para mejorar la experiencia de usuario y facilitar futuras modificaciones del código.
    \item \textbf{Sprint 9 (julio 2026 - septiembre 2026)} Realización de pruebas experimentales para validar el correcto funcionamiento de la aplicación web y su contribución al desarrollo del pensamiento computacional. Análisis de los resultados obtenidos y redactar la memoria del proyecto.
    \item \textbf{Sprint 10 (septiembre 2026)} Entrevista a Jens Mönig y experimento con profesores de secundaria. Diseño de experimento con alumnos en el aula. Incorporación de los resultados y experimento diseñado a la memoria del proyecto y entrega de la misma.
\end{itemize}

La consecución de dichos Sprints en el orden descrito se puede revisar en el repositorio de GitHub siguiente, donde, además, se encuentra el código fuente de la aplicación web implementada: \url{https://github.com/miribuenot/Dr_Snap/tree/main}

Asimismo, se muestra en la Figura \ref{fig:gantt_sprints_colores} un diagrama de Gantt con la planificación temporal de los Sprints y las tareas que se han llevado a cabo en cada uno de ellos.

\begin{figure}[htbp]
    \centering
    \resizebox{\textwidth}{!}{ % Ajusta el diagrama al ancho de la página
    \begin{ganttchart}[
        hgrid,
        vgrid,
        x unit=1.4cm,
        y unit title=0.6cm,
        y unit chart=0.7cm,
        title label font=\bfseries\footnotesize,
        bar label font=\footnotesize,
        bar height=0.6,
    ]{1}{11}

    % Títulos de años y meses
    \gantttitle{2025}{2} \gantttitle{2026}{9} \\
    \gantttitle{Nov}{1} 
    \gantttitle{Dic}{1} 
    \gantttitle{Ene}{1} 
    \gantttitle{Feb}{1} 
    \gantttitle{Mar}{1} 
    \gantttitle{Abr}{1} 
    \gantttitle{May}{1} 
    \gantttitle{Jun}{1} 
    \gantttitle{Jul}{1} 
    \gantttitle{Ago}{1} 
    \gantttitle{Sep}{1} \\

    % Sprints con colores personalizados
    \ganttbar[bar/.style={fill=red!30, draw=red!80, rounded corners=2pt}]{Sprint 1: Análisis Inicial}{1}{2} \\
    \ganttbar[bar/.style={fill=orange!30, draw=orange!80, rounded corners=2pt}]{Sprint 2: API \& Traducción}{2}{3} \\
    \ganttbar[bar/.style={fill=yellow!40, draw=yellow!90, rounded corners=2pt}]{Sprint 3: Batch Mode \& Parseo}{3}{4} \\
    \ganttbar[bar/.style={fill=green!30, draw=green!80, rounded corners=2pt}]{Sprint 4: Scripts Duplicados}{4}{4.45} \\
    \ganttbar[bar/.style={fill=teal!30, draw=teal!80, rounded corners=2pt}]{Sprint 5: Logo \& Interfaz}{4.55}{5} \\
    \ganttbar[bar/.style={fill=cyan!30, draw=cyan!80, rounded corners=2pt}]{Sprint 6: Ciberseguridad \& Versión Beta}{5}{6} \\
    \ganttbar[bar/.style={fill=blue!30, draw=blue!80, rounded corners=2pt}]{Sprint 7: Despliegue}{6}{7} \\
    \ganttbar[bar/.style={fill=violet!30, draw=violet!80, rounded corners=2pt}]{Sprint 8: Refactorización}{7}{9} \\
    \ganttbar[bar/.style={fill=magenta!30, draw=magenta!80, rounded corners=2pt}]{Sprint 9: Pruebas y Memoria}{9}{11} \\
    \ganttbar[bar/.style={fill=gray!30, draw=gray!80, rounded corners=2pt}]{Sprint 10: Entrevistas y Cierre}{11}{11}

    \end{ganttchart}
    }
    \caption{Diagrama de Gantt de los Sprints del proyecto.}
    \label{fig:gantt_sprints_colores}
\end{figure}
